\documentclass[11pt, a4paper]{article}

\usepackage[T1]{fontenc} 
\usepackage{cmbright}

% --- Essential Math Packages ---
\usepackage{amsmath}    % Core math environments
\usepackage{amssymb}    % Math symbols
\usepackage{mathtools}  % Improvements to amsmath (fixes spacing)
\usepackage{bm}         % Bold math (better than \mathbf for greek letters)

% --- Formatting & Utilities ---
\usepackage[margin=1in]{geometry} % Standard margins
\usepackage{microtype}            % Subliminally improves spacing/kerning
\usepackage{hyperref}             % Clickable links in PDF
\usepackage{cleveref}             % Smart referencing (use \cref instead of \ref)
\usepackage{url}

% --- Settings ---
\hypersetup{
    colorlinks=true,
    linkcolor=blue,
    citecolor=green,
    urlcolor=cyan
}

\title{Note on reverse-mode linear solver gradients}
\author{Lu Hong}
\date{\today}

\begin{document}

\maketitle

\paragraph{Adjoint method}
Let $A\in\mathbb R^{N\times N}$ be an invertible matrix and $b\in\mathbb R^{N}$ such that there is a unique $x\in\mathbb R^N$ that satisfies

\begin{equation}
    Ax=b \label{linear_sys}
\end{equation}

Let $L:\mathbb R^N\to \mathbb R$ be a loss function of $x$. The goal of this note is to derive an expression for the gradient of $L$ with respect to $A$ and $b$.

To find these gradients, we use the adjoint method and start by defining the Lagrangian
\begin{equation}
    F(x, A, b, \lambda) = L(x) - \left<\lambda, Ax-b\right> \label{lagrangian}
\end{equation}
where $\lambda\in\mathbb R^N$ is a free vector. By this definition, when $x$ is a valid solution to \eqref{linear_sys}, the second term in \eqref{lagrangian} vanishes, which implies that $F=L$; a consequence of this equality is that the total derivatives of $F$ and $L$ with respect to $A$ and $b$ are also identical.

Therefore, if $x$ is a solution to \eqref{linear_sys}, we can compute the total derivative (i.e., gradient) of $L$ with respect to $A$ as
\begin{equation}
    \frac{dL}{dA} = \frac{dF}{dA} = \frac{\partial F}{\partial x}\frac{\partial x}{\partial A} + \frac{\partial F}{\partial A} \label{partial_A}
\end{equation}
where the second equality follows from the fact that $F$ depends on A both directly and indirectly via $x(A, b)$. Note that, to compute $\partial x/\partial A$ would require tracking gradients through the linear solver, which is expensive. Since $\lambda$ is a free vector, we can choose a $\lambda$ such that the first term in \eqref{partial_A} vanishes and \eqref{partial_A} reduces to
\begin{equation}
    \frac{dL}{dA} = \frac{\partial F}{\partial A} \label{partial_A_simplified}
\end{equation}
This is equivalent to the requirement that
\begin{align}
    \frac{\partial F}{\partial x} & = 0 \\
    \frac{\partial L}{\partial x} - \frac{\partial}{\partial x}\left<\lambda, Ax\right> & = 0 \\
    \frac{\partial}{\partial x}\left<A^T\lambda, x\right> & = \frac{\partial L}{\partial x} \\ 
    A^T\lambda & = \frac{\partial L}{\partial x} \label{lambda_linear_sys}
\end{align}
Note that the adjoint property of $A$ is the key step in enabling the simplification of the derivative of the inner product $\left<\lambda, Ax\right>$ with respect to $x$ by moving the action of $A$ to the other side of the inner product (hence the name ``adjoint method'').

Taken together, once we solve \eqref{lambda_linear_sys} for $\lambda$, we can use \eqref{partial_A_simplified} to find the total derivative of $L$ with respect to $A$ as
\begin{equation}
    \frac{dL}{dA} = \frac{\partial F}{\partial A} = \frac{\partial L}{\partial A} + \frac{\partial}{\partial A}\left<-\lambda, Ax-b\right> = -\lambda x^T
\end{equation}
note that $\partial L/\partial A$ is zero since $L(x)$ has no explicit dependency on $A$. Repeating the same argument for $b$, we find that the same $\lambda$ can be used to find the total derivative of $L$ with respect to $b$ as
\begin{equation}
    \frac{dL}{db} = \frac{\partial F}{\partial b} = \frac{\partial}{\partial b}\left<-\lambda, Ax-b\right> = \lambda
\end{equation}

\paragraph{Channel dimension}
Sometimes, it is convenient to define a matrix $B\in \mathbb R^{N\times C}$ with a batch/channel dimension and solve the linear system
\begin{equation}
    AX=B
\end{equation}
for the unknown matrix $X\in\mathbb R^{N\times C}$, which is akin to solving $C$ linear systems of the form \eqref{linear_sys} in parallel. In this case, the $\lambda$ vector is replaced with a matrix $\Lambda$, which results in the Lagrangian
\begin{equation}
    F(X, A, B, \Lambda) = L(X) - \left<\Lambda, AX-B\right>_F
\end{equation}
where $\left<\cdot, \cdot\right>_F$ is the Frobenius inner product between matrices defined as
\begin{equation}
    \left<A, B\right>_F = \text{Tr}(A^T B)
\end{equation}

Retracing the previous derivation using the new inner product definition, we find that $\Lambda$ is the solution to the linear system
\begin{equation}
    A^T\Lambda = \frac{\partial L}{\partial X}
\end{equation}
and the total derivative of $L$ with respect to $A$ is
\begin{equation}
    \frac{dL}{dA} = -\frac{\partial}{\partial A}\left<\Lambda, AX\right>_F = -\frac{\partial}{\partial A}\text{Tr}\left(\Lambda^TAX\right) \overset{1}{=} -\frac{\partial}{\partial A}\text{Tr}\left(X\Lambda^TA\right) \overset{2}{=} -\Lambda X^T \label{grad_A_contract}
\end{equation}
Here, (1) follows from the fact that the trace of matrix multiplication is invariant to cyclic permutations of the matrices, and (2) follows from the fact that
\begin{equation}
    \frac{\partial}{\partial A}\text{Tr}(BA) = B^T
\end{equation}
Note that, if we write $\Lambda = [\lambda_1, \cdots, \lambda_C]$ and $X = [x_1, \cdots, x_C]$ as a collection of channel column vectors, then \eqref{grad_A_contract} is equivalent to
\begin{equation}
    \frac{dL}{dA} = -\sum_{c=1}^C \lambda_c x_c^T
\end{equation}
that is, $dL/dA$ can be computed independently as outer products for each channel and summed together.

Lastly, the total derivative of $L$ with respect to $B$ is
\begin{equation}
    \frac{dL}{dB} = \frac{\partial}{\partial B}\left<\Lambda, B\right>_F = \Lambda
\end{equation}

\paragraph{Uniqueness condition}
Note that, since the adjoint method depends on $x(A, b)$ being a well-defined function of $A$ and $b$, it breaks down if the solution to \eqref{linear_sys} is not unique. In such cases, some methods will fail to find a solution (e.g., those based on LU factorization), in which case the gradients are undefined. However, some linear solvers (e.g., those computing the Moore-Penrose pseudoinverse via SVD or LSQR) enforce a minimum-norm condition if \eqref{linear_sys} has non-unique solutions (i.e., the uniqueness condition is guaranteed by finding $x$ with the smallest $\|x\|^2$); in such cases, the adjoint method presented above gives incorrect gradients, because it does not account for the minimum-norm constraint. In cases where \eqref{linear_sys} is close to singular, the gradients given by the adjoint method are still correct, but the solution to \eqref{lambda_linear_sys} will be numerically unstable since the condition number of $A$ is large.

\end{document}